%======================================================================
%  Course overview -- a single page, neither chapter nor section
%======================================================================
\cleardoublepage
\phantomsection
\addcontentsline{toc}{chapter}{Course overview}
\markboth{Course overview}{Course overview}
\thispagestyle{plain}

\noindent{\color{gray!70!black}\rule{\textwidth}{0.6pt}}\par\vspace{6pt}\noindent
{\sffamily\bfseries\Huge Course overview\par}
\vspace{18pt}

\paragraph{Stability in mechanics.}
A structure is loaded slowly. For a long time it deforms a little more
under each increment. Then, at a critical load, a new shape becomes
possible and the structure takes it: a column bows, a plate wrinkles, a
cylinder collapses into diamonds. The course explains when this happens,
what the structure does next, and why some structures fail far below the
load predicted by the theory. On 29 August 1907 the south arm of the Quebec Bridge
fell during construction because a compressed chord buckled. The same mathematics serves in the opposite
direction: deployable and morphing structures use instability on purpose.

\paragraph{Roadmap.}
The course goes up in dimension: from one unknown to finitely many, then
to rods, plates, shells and three-dimensional bodies.
\begin{enumerate}[itemsep=1pt]
  \item Differential equations and bifurcation (Chapter~\ref{ch:st}).
  \item Discrete systems: energy, Lagrangian mechanics and Koiter's analysis (Chapter~\ref{ch:ld}).
  \item From continuum to discrete: numerical methods (Chapter~\ref{ch:nc}).
  \item Rods: the elastica and beyond (Chapter~\ref{ch:rd}).
  \item Rods II: thin-walled members, cruciform sections and residual stresses (Chapter~\ref{ch:tw}).
  \item Plates (Chapter~\ref{ch:pl}).
  \item Cylindrical shells, imperfection sensitivity and rocket fuselages (Chapter~\ref{ch:cy}).
  \item Shells in service: nozzles, prestress and deployable structures (Chapter~\ref{ch:sh}).
  \item Three-dimensional bodies (Chapter~\ref{ch:so}).
\end{enumerate}

\paragraph{One method, four tools.}
Every chapter follows the same path. A closed-form solution comes first,
on the simplest model that shows the effect. Mathematica carries the
asymptotic expansions. A short Python code follows the branches by
continuation. FEniCS and Cast3M solve the full problem by finite elements,
and are trusted only after they reproduce the closed form.

\paragraph{What you should be able to do at the end.}
\begin{itemize}[itemsep=1pt]
  \item find critical points and decide their stability from the energy;
  \item tell a stable post-buckling path from an imperfection-sensitive one;
  \item follow equilibrium branches through limit and bifurcation points;
  \item set up and check a buckling computation by finite elements;
  \item explain the knockdown factors used to design thin shells.
\end{itemize}

\paragraph{Practical matters.}
Nine sessions of two hours of lecture and two hours of exercises.
Prerequisites: the variational toolkit and the finite element method of
the notes \emph{Nonlinear Problems in Mechanics}.
The line of the course is structural engineering, with some physical
examples. Main texts: Thompson \& Hunt, Ba\v{z}ant \& Cedolin, Brush \&
Almroth, Koiter, and Guckenheimer \& Holmes for the dynamical systems.
Design practice: Bushnell's survey of shell buckling and the NASA
monographs SP-8007. \draftnote{Physical examples (Audoly \& Pomeau, Reis,
Hutchinson, Bertoldi) and the final list to be settled.}

\paragraph{Assessment.}
\draftnote{Proposal: written exam and a project in pairs, about one hour per session from Lecture~3 on,
on one real structure (fuselage panel, nozzle, deployable boom, cruciform
column), combining a closed form, a continuation and a finite element run.}

\cleardoublepage
